%! Introducing Vectors — Unit 1, Lesson 1.0 — Student Packet Cover Sheet
\documentclass[10pt]{article}
\usepackage{linalg-article}
\usepackage{linalg-boxes}
\usepackage{ltablex}
\keepXColumns

\begin{document}

% Full-bleed burgundy banner
\begin{tikzpicture}[remember picture, overlay]
  \fill[burgundy] (current page.north west)
    rectangle ([yshift=-0.9in]current page.north east);
\end{tikzpicture}
\vspace{-0.8in}

\begin{center}
  {\color{white}\textbf{\LARGE Linear Algebra}}\\[4pt]
  {\color{blushmid}\large\textbf{Unit 1 \quad Vectors and Matrices}}\\[6pt]
  {\color{white}\textbf{\large Lesson 1.0 \quad Introducing Vectors --- Vocabulary and Length}}
\end{center}

\vspace{0.22in}
\namedateperiod
\vspace{0.18in}

\begin{learningtargetbox}
\small
\begin{enumerate}[leftmargin=1.5em, itemsep=5pt]
  \item \ldots{} describe a vector by its \textbf{components} and as an \textbf{arrow} in the plane.
  \item \ldots{} \textbf{add} vectors and multiply a vector by a \textbf{scalar}, and picture what happens to the arrow.
  \item \ldots{} find the \textbf{length (magnitude)} of a vector and explain what it measures.
\end{enumerate}
\end{learningtargetbox}

\vspace{0.14in}

\begin{tocbox}
\small
\renewcommand{\arraystretch}{1.5}
\begin{tabularx}{\linewidth}{@{} c l X r @{}}
\rowcolor{blush}
\textbf{\#} & \textbf{Component} & \textbf{Description} & \textbf{Score} \\
\midrule
1 & Warm-Up        & Spiral review: signed numbers, Pythagoras, coordinates & \blank{1.2cm} \\
2 & Guided Notes   & What a vector is; adding, scaling, and length            & \blank{1.2cm} \\
3 & Group Activity & Drone Delivery --- combining hops and measuring distance  & \blank{1.2cm} \\
4 & Exit Ticket    & Quick check on components and magnitude                   & \blank{1.2cm} \\
5 & Homework       & Practice with vectors, length, and a data-record extension & \blank{1.2cm} \\
\midrule
  & \multicolumn{2}{r}{\textbf{Total}} & \blank{1.2cm} \\
\end{tabularx}
\end{tocbox}

\vspace{0.10in}

\begin{remindbox}
\small A \textbf{vector} is an ordered list of numbers --- two numbers in the plane. Picture it
as an \textbf{arrow} from the origin. Its \textbf{length} is
$\|\mathbf{v}\|=\sqrt{v_1^{\,2}+v_2^{\,2}}$ (the Pythagorean theorem), and it is never negative.
\end{remindbox}

\end{document}
